\subsubsection{7. Endpoint coupling, hybrid identity, and conditional moment route}\label{endpoint-coupling-hybrid-identity-and-conditional-moment-route}

\paragraph{7.1 Joint endpoint functional TV}\label{joint-endpoint-functional-tv}

Let \((z,z^k)\) be coupled product inputs and let \(x=F(z),y=F(z^k)\), with truncations \(x^R=F^R(z),y^R=F^R(z^k)\). Suppose \(A_S\) is an event on which both truncated endpoints agree with the untruncated ones on \(S\), with
\[\Pr[A_S^c]\le\eta(M_S,R)=2C_0M_Se^{-c_0\delta R}+2L2^{-L}.\]
Then, for every \([0,1]\)-valued test functional \(\Phi\),
\[|\mathbb E\Phi(x_S^R,y_S^R)-\mathbb E\Phi(x_S,y_S)|
\le\Pr[A_S^c]\le\eta(M_S,R).\]
The proof is simply that the difference vanishes on \(A_S\) and is at most \(1\) off it. An ordered path functional defined from \textbf{both} endpoints is an admissible \(\Phi\) if it depends only on the displayed coordinates.

This proves a joint-law TV comparison only. It does not show that an intermediate path value is dominated by the original squared gradient or by \(\operatorname{Inf}_i^\mu\); that separate missing step cannot be inferred from the TV estimate.

\paragraph{7.2 Exact hybrid/Taylor identity}\label{exact-hybridtaylor-identity}

Let \(X\sim\mu\), \(Y\sim\nu\) independently, and define
\[Z^{(i)}=(Y_1,\ldots,Y_i,X_{i+1},\ldots,X_n),\qquad
W^{(i)}=(Y_1,\ldots,Y_{i-1},X_{i+1},\ldots,X_n).\]
For any \(h\), \(\mathbb E_\mu h-\mathbb E_\nu h=\sum_i\mathbb E[h(Z^{(i-1)})-h(Z^{(i)})]\) exactly. Multilinearity gives
\[f(Z^{(i-1)})-f(Z^{(i)})
=(X_i-Y_i)\,\partial_ip(W^{(i)})/\sqrt v.\]
For \(\psi\in C^3\), exact Taylor expansion gives
\[\begin{aligned}
\psi(f^{(i-1)})-\psi(f^{(i)})
={}&\psi'(f^{(i)})(X_i-Y_i)\partial_ip/\sqrt v\\
&+\tfrac12\psi''(f^{(i)})(X_i-Y_i)^2(\partial_ip)^2/v+R_i,
\end{aligned}\]
with
\[|R_i|\le\frac16\|\psi^{(3)}\|_\infty
|X_i-Y_i|^3|\partial_ip|^3/v^{3/2}.\]
These are identities only. The first-order covariance bound and any subsequent exponential-tail comparison remain open. A smooth \(C^3\) error polynomial in \(t\) does not prove a lower tail of size \(e^{-\Theta(t^2)}\).

\paragraph{7.3 Conditional HC+CU route}\label{conditional-hccu-route}

Assume, for standardized \(f\) with mean \(0\), variance \(1\), and \(I=\max_i\operatorname{Inf}_i(f)\), both
\[\|f\|_q\le(K_0q)^{K_1d}\quad\text{(HC)},\]
and, for \(q\ge3\),
\[|\kappa_q(f)|\le(C_1dq)^{C_2dq}I^{q/2-1}q!\quad\text{(CU)}.\]
After absorbing model-dependent constants, \(\log M(s)=\sum_q\kappa_qs^q/q!\) gives, on the stated geometrically small interval,
\[\log M(s)=\frac{s^2}{2}+O(s^3\sqrt I).\]
Tilting by \(dP_s=e^{sf}dP/M(s)\) differentiates the cumulant function once and twice. The tilted mean is \(s+O(s^2\sqrt I)\) and the tilted variance is \(1+O(s\sqrt I)\). Choosing \(s=t+O(1)\) and using Chebyshev on the tilted law yields, in the admissible interval,
\[\Pr[f>t]\ge c\exp(-t^2/2+O(t^3\sqrt I)).\]
Apply the same argument to \(-f\). Thus, whenever \(t^3\sqrt I=O(1)\), the probability is positive with an \(n\)-independent lower bound. If \(I\le c t^{-6}\) with suitable \(c\), this condition holds.

For bounded \(t\), HC and \(Ef^4\le K^4\) give
\[\Pr(|f|>1/2)\ge\frac{(1-1/4)^2}{K^4}=\frac9{16K^4},\]
by splitting \(Ef^2\) at level \(1/2\) and applying Cauchy--Schwarz on the large part.

Finally, for \(p\in[0,1]\) and \(f=(p-\mathbb Ep)/\sqrt v\), boundedness makes \(\Pr(|f|>2/\sqrt v)=0\). Taking \(t=2/\sqrt v\) in the contrapositive gives, polynomially,
\[I\gtrsim t^{-6}\asymp v^3,
\qquad \max_i\operatorname{Inf}_i(p)=vI\gtrsim v^4.\]
This is a complete \textbf{conditional} proof only. HC and CU are not established under general Dobrushin, and the moment argument does not prove them.
